\section*{Chapter \ref{ch:s2:flow-market-making-at-a-bank} --- Flow Market Making at a Bank}

\begin{solution}{exo:s2:flow-market-making-at-a-bank:1}
$(1.2 - 0.83) \times \$100 = \$37$ per trade, before hedging costs.
\end{solution}

\begin{solution}{exo:s2:flow-market-making-at-a-bank:2}
The skew is $0.05 \times 4 = 0.2$ basis points against the long position: a buying client, who reduces it, pays $1.2 - 0.2 = 1.0$; a selling client,
who adds to it, pays $1.2 + 0.2 = 1.4$.
\end{solution}

\begin{solution}{exo:s2:flow-market-making-at-a-bank:3}
From $0.8 = 1/(1 + e^{(1.2 - c)/0.25})$: $c = 1.2 + 0.25 \ln(0.8/0.2) = 1.55$ basis points.
\end{solution}

\begin{solution}{exo:s2:flow-market-making-at-a-bank:4}
A single price must serve clients with different mark-outs and different competing quotes. A low price wins the informed clients, whose trades lose
money; a high one loses captive clients' volume along with the informed. Pricing by client charges captive clients what they will pay and prices
informed clients out, which no single number can do.
\end{solution}

\begin{solution}{exo:s2:flow-market-making-at-a-bank:5}
A small skew costs little capture and saves hedging and inventory marks: clients trade more in the direction that reduces the position. A large skew
gives away capture on many trades to save hedging costs that are already nearly zero: \$7.63 million at 0.05, \$7.08 million at 0.2.
\end{solution}

\begin{solution}{exo:s2:flow-market-making-at-a-bank:6}
A dealer can skew its prices to attract offsetting flow; price-sensitive clients respond to the skew, so the position is offset sooner and at a smaller
concession, as Butz and Oomen found.
\end{solution}

\begin{solution}{exo:s2:flow-market-making-at-a-bank:7}
Trades rise to 83\,027, marks fall to $-\$1.65$ million and the \omterm{def:s2:flow-market-making-at-a-bank:franchise}{franchise P\&L} to \$4.60 million. The informed clients now trade 3\,259 times at a
net $-1.72$ basis points each, and even the moderate ones lose the desk money ($-0.07$).
\end{solution}

\begin{solution}{exo:s2:flow-market-making-at-a-bank:8}
Volume from informed clients is the desk's loss: on the flat price their trades netted $-1.18$ basis points each, and tighter prices would win more of
them. Their mark-outs, not their volume, should set their prices, and the \omterm{def:s2:flow-market-making-at-a-bank:franchise}{franchise P\&L}, not volume, measures the desk.
\end{solution}

\begin{solution}{pb:s2:flow-market-making-at-a-bank:1}
\begin{enumerate}
\item The profit from serving clients: spread capture, inventory marks (where adverse selection shows) and hedging costs.
\item 60 uninformed, 20 moderately informed (0.8 basis point drift) and 10 informed (2.5) clients, each with its own competing quote; the drift follows
the client's trade whoever fills it, so the mid path is common to all policies.
\item Mark-outs of $-0.01$, 0.84 and 2.41 basis points and hit ratios of 80\%, 43\% and 8\%; the inferred competing quotes correlate 0.998 with the
true ones.
\item Price signatures, the average price path around executions, tell internalising liquidity providers from externalising ones.
\item Pricing each client by its mark-out and its sensitivity to price.
\item The half-spread that maximises (half-spread less mark-out) times the win probability, with the competing quote inferred from the hit ratio.
\item Flat 1.0, 1.2, 1.4, 1.6: \$5.90, 6.57, 6.63, 6.23 million; priced by client \$7.60 million, with the skew \$7.63 million.
\item Uninformed: capture 1.20 to 1.39, net 1.20 to 1.40 basis points; moderate: 1.20 to 1.25, net 0.37 to 0.44; informed: 938 trades at $-1.18$ to
10 trades.
\item A shift of the quote against the inventory or axes, to attract trades that reduce the position.
\item At 0.05: P\&L \$7.63 million, average inventory 2.18 units instead of 2.75, internalisation 96.8\% instead of 90.8\%; at 0.2: \$7.08 million,
1.30 units, all internalised.
\item Internalisation horizons of several minutes in liquid currencies and tens of minutes in emerging markets for a representative large dealer;
lower costs for dealers holding more risk and facing price-sensitive clients.
\item The dealer's option to reject or re-price a deal request after latency; its design sets the client's execution risk but need not change its
effective costs.
\item Priced by client, the desk captures 1.39 basis points per trade from uninformed clients and 1.25 from moderately informed ones, prices the
informed out, and internalises 90.8\% of its volume, 96.8\% with a 0.05 skew.
\item That the win curve's shape is known (logistic with a known width): the desk inferred competing quotes almost exactly.
\item It wins the informed clients' trades and loses on each.
\item Vary quotes deliberately (small randomised offsets) and fit hit ratios against them, by client segment and size.
\item Re-price it: widen its spread or reduce its size, and check whether its behaviour changed or its mark-out estimate was noisy.
\item Internalisation before hedging, which holds the risk while waiting for offsetting flow.
\item Chapter~15 trades against predictable flows from outside; this chapter is the desk that receives flows and prices them.
\item Immediacy, priced to each client's information and alternatives.
\end{enumerate}
\end{solution}

\begin{solution}{iq:s2:flow-market-making-at-a-bank:1}
The mid's move after a trade, measured in the client's favour, over one or several horizons. Averaged over a client's trades it is what the desk loses
to that client's information; add it to the target margin to set the client's spread.
\end{solution}

\begin{solution}{iq:s2:flow-market-making-at-a-bank:2}
Use all requests, won and lost: the hit ratio against the price quoted. Randomise quotes slightly so the curve is identified, fit a model of the win
probability by client segment, size and market conditions, and correct for the fact that quotes were set with information about the client.
\end{solution}

\begin{solution}{iq:s2:flow-market-making-at-a-bank:3}
A skewed two-way price: a tighter offer to sell down the position and a wider bid, sized to the client's mark-out, the position and the cost of
hedging; not a one-sided refusal, which damages the franchise.
\end{solution}

\begin{solution}{iq:s2:flow-market-making-at-a-bank:4}
From the expected time to offsetting flow and the volatility over that time: the limit's loss at a stress move must fit the desk's risk budget; test
it against one-way markets, when offsetting flow does not come.
\end{solution}

\begin{solution}{iq:s2:flow-market-making-at-a-bank:5}
Subscribe to trades and to a reference mid; for each trade schedule mark-out computations at fixed horizons; aggregate by client with decaying
weights; publish estimates with confidence to the pricing engine; keep the raw data for research and audit.
\end{solution}

\begin{solution}{iq:s2:flow-market-making-at-a-bank:6}
Maximise $(h - m)p(h)$: $p + (h - m)p' = 0$, and for the logistic $p' = -p(1-p)/w$, so $p - (h - m)p(1-p)/w = 0$ and $h - m = w/(1 - p(h))$. The
margin over the mark-out rises as the win probability rises: captive clients pay more.
\end{solution}
